\cotefile{157-5}
\scanpage{69}\mspage{56}
% archive sheet number: 68

$f\in W(I^{\unsure{o|\circ}})$, alors $i\in W(I^o)$ donc $i\in W(I^o)\cap C_{I^o}=(TC)_{I^o}$.
Au total, \deleted{\unsure{tout}} $f\in W(I^o)_c$ est composé $pi$\inserted{ dans $M(I^o)_c$},
où $i\in(TC)_{I^o}$ et $p$ est inverse à gauche
d'un $j\in TC$. Comme $u_!$ transforme
$(TC)_{I^o}$ en $(TC)_{I'^o}$, il transforme $f$
en $f'=p'i'$, avec $i'\in(TC)_{I'^o}\subset W_{I'^o}$,
et $p'$ inverse à gauche d'un $j'=u_!(j)$
qui est \unsure{également} dans $W_{I'^o}$. Mais alors
$\deleted{\unsure{j'}}\unsure{p'|j'}\in W(I'^o)$, donc aussi $f'=p'i'\in W_{I'^o}$,
ce qui prouve que $u_!(W(I^o)_c)\subset W(I'^o)$.

La dernière assertion provient du \deleted{\unsure{fait}}\inserted{lemme}
suivant, appliqué à $(M(I^o),W(I^o),C_{I^o})$ \unsure{ald} $(M,W,C)$\unsure{.}

\uline{Lemme} \uline{3.4}. Soit $(M,W,C)$ une catégorie
à cofibrations. Alors l'inclusion $M_c\subset M$
induit une équivalence \deleted{\unsure{$\mathrm{Ho}_W$\illegible{1}}}
\[
M_c W_c^{-1}\overset{\approx}{\unsure{\to|\text{---}|\leftrightarrow}}MW^{-1}=\mathrm{Ho}_W\,.
\]

Nous donnerons la démonstration (par
\unsure{A la J|A le J|A et J|A 6 J|A \& J}) ultérieurement\unsure{,|.}

Cela achève la démonstration de
3.3.

3.5. Il se pose la question si \deleted{\unsure{le}} le foncteur
$u_!^W$ de (3.3.1) est bien un adjoint à g.
\scanpage{70}\mspage{57}
% archive sheet number: 69
% Paragraph continues from the preceding sheet.
du foncteur en sens inverse. $u_W^*$, induit
par $u^*$. Notons que c'est le cas si on
\deleted{suppose que} $u^*$ transforme $M(I'^{\unsure{o|\circ}})_c$ en
$M(I^o)_c$, et quand on suppose que
$I'$ \inserted{est également artinien, et} satisfait \unsure{également} à une des hypothèses
a) b) du th. 3.1. Car alors $(M(I'^o),W(I'^o),C_{I'^o})$
est une catégorie à cofibrations par le \unsure{th.|Th.}
3.1, on peut \unsure{leur|lui} appliquer le lemme
précédent, donc $u_W^*$ et $u_!^W$ se \unsure{réduisent|déduisent} à
partir d'un couple de foncteurs

\diagram{\begin{tikzcd}[ampersand replacement=\&,column sep=large]
M(I^o)_c \arrow[r,dash,shift left=1.4ex,"u_!^c", "\unsure{\text{---}|\to}"{description,near end,above=1pt,font=\tiny}] \& M(I'^o)_c\,, \arrow[l,shift left=1.4ex,"u_c^*"]
\end{tikzcd}}

lesquels sont visiblement adjoints l'un
de l'autre (puisque induits par un
couple de foncteurs adj. $u_!,u^*$), et
de plus compatibles aux localiseurs.
Donc on peut appliquer le lemme des
foncteurs \unsure{adjoint|adjoints} pour les localisés. Pour
résumer :

\uline{Corollaire} 3.5. Supposons que $I,I'$ sont
\deleted{\unsure{loc}} artiniens, et qu'on soit dans l'un
des deux cas \unsure{suivant|suivants}

a) \deleted{$I,I'$ sont \unsure{no}} les espaces $\mathrm{Esp}(I),\mathrm{Esp}(I')$
sont noethériens (p. ex. $I,I'$ finis).
\scanpage{71}\mspage{58}
% archive sheet number: 70

b) $C$ et $TC=C\cap W$ sont stables par composition
transfinie.

Supposons de plus que $u^*$ applique $M(I'^{\unsure{o|\circ}})_c$
dans $M(I^o)_c$, ce qui est le cas notamment
si $u$ est loc. \deleted{\unsure{filtr.}} cofiltrant i.e. les
\deleted{\unsure{cat.}}\inserted{ordonnés} localisés $\mathord{}_{i'}\backslash I$ sont loc. cofiltrants. \deleted{\unsure{à g.}}
Alors le foncteur $u_!^W$ construit en (\unsure{3.3.1|3.2.1}).
est adjoint à g. de $u_W^*$.

\uline{Remarque} 3.6. Quand on a un couple
de foncteurs adjoints

\diagram{\begin{tikzcd}[ampersand replacement=\&,column sep=large]
C \arrow[r,dash,shift left=1ex,"u_!", "\unsure{\text{---}|\to}"{description,near end,above=1pt,font=\tiny}] \& C', \arrow[l,shift left=1ex,"u^*"]
\end{tikzcd}}

avec des localiseurs $W,W'$ de $C,C'$,
$u^*$ compatible aux localiseurs, et si de
plus on a une sous-catégorie \inserted{pleine} $C_0$ de $C$
telle que a) $\deleted{\unsure{C}}u_!|C_0$ transforme flèches
de $W$ en flèches de $W'$ et b) l'inclusion
$C_0\hookrightarrow C$ induit une \uline{équivalence}
$C_0W_0^{-1}\overset{\approx}{\text{---}}CW^{-1}$ (où $W_0=W\cap\mathrm{Fl}(C_0)$),
on trouve deux foncteurs canoniques

\diagram{\begin{tikzcd}[ampersand replacement=\&,column sep=large]
C W^{-1} \arrow[r,dash,shift left=1.2ex,"u_!^W"] \& C' W'^{-1}, \arrow[l,shift left=1.2ex,"u_W^*"]
\end{tikzcd}}

où $u_!^W$ est défini \inserted{seulement} à isom \unsure{can$^{\mathrm{l}}$|can$'$} près.
\scanpage{72}\mspage{59}
% archive sheet number: 71

Contrairement à ce que prétend Anderson,
il n'est pas vrai en général que ces
deux foncteurs soient adjoints, sans hypothèses
supplémentaires. Voici un
\unsure{contreexemple|contre-exemple} : on prend $C=C'$, $u_!=u^*=\mathrm{id}_C$,
$W'=\mathrm{Is}_{\deleted{\unsure{e|C}}}(C)$, $W=\mathrm{Fl}(C)$, donc \deleted{\unsure{le}} $u_W^{\unsure{'|1|*}}$ est
le foncteur canonique

\diagram{\begin{tikzcd}[ampersand replacement=\&,column sep=large]
\underbrace{\Pi C}_{\substack{\text{groupoïde}\\\text{fondamental}}} \& C\, . \arrow[l]
\end{tikzcd}}

Supposons que $C$ soit 1-connexe, \deleted{\unsure{donc}} i.e
$\deleted{\unsure{C}}\Pi C$ soit équiv. à $e$. Soit $e\in\mathrm{Ob}\,C$
tel que \deleted{le \unsure{groupe} enveloppant du
monoïde} $\mathrm{End}_C(e)$ soit le groupe unité,
\deleted{p. ex. tel que $\mathrm{End}_C(e)=\{\mathrm{id}_e\}$,} \deleted{Alors} et
prenons $C_0=\{e\}$. Alors l'hypothèse
sur $C_0$ est satisfaite. Mais si $u_!^W$
était adjoint à gauche \inserted{de $u_W^*$}, $e$ serait objet
initial de $C$ ; \deleted{\unsure{donc} des objets rigides}
par exemple si $C$ a un objet final $e$
(ce qui implique que $C$ est 1-connexe et
$e$ rigide), cet objet final devrait
être initial si Anderson avait raison avec
``it is clear that \ldots''.
\scanpage{73}\mspage{60}
% archive sheet number: 72

Néanmoins :

\uline{Théorème} 3.7. Soit $(M,W,C)$ une catégorie à
cofibrations \inserted{stable par $\underrightarrow{\lim}$ finies}, et $u:I\unsure{\text{---}|\to}I'$ une application
croissante entre ens. ordonnés artiniens.
On suppose qu'on est sous l'une des
hypothèses suivantes

a) $\mathrm{Esp}(I),\mathrm{Esp}(I')$ sont noethériens, p. ex.
$I,I'$ finis

b) $C$ et $TC\overset{\mathrm{déf}}{=}C\cap W$ sont stables par composition
transfinie, et $M$ stable par \inserted{petites} $\underrightarrow{\lim}$ filtrantes.

Alors le foncteur $u_!^W:\mathrm{HO}_W(I^{\unsure{o|\circ}})\mathrel{\text{---}}\mathrm{HO}_W(I'^o)$
construit dans (3.3.1) est adjoint à gauche
du foncteur $u_W^*$ en sens opposé.

\uline{Dém.} On utilise le diagramme bien
connu

\diagram{\begin{tikzcd}[ampersand replacement=\&,column sep=huge,row sep=large]
\& I \arrow[d,"\partial_1"'] \arrow[dd,bend left=50,"w"] \\
I \arrow[r,"{\partial_0\ (\mathrm{ouv.})}","\unsure{\to|\leftrightarrow}"{description,very near start,above=1pt,font=\tiny}] \arrow[d,"u"'] \arrow[dr,phantom,"\text{\unsure{cocart}}"{description}] \& I\times\unsure{\mathord{\tikz[baseline=0pt,x=1em,y=1em]{\draw[line width=.045em,line join=round] (0,.02)--(.42,.68) (.17,.02)--(.52,.62)--(.83,.02)--cycle;}}|\Delta}^{1} \arrow[l,bend left=18,dashed,"\mathrm{pr}_1"] \arrow[d] \\
I' \arrow[r,"\text{imm ouv.}"] \& C(u) \arrow[l,bend left=18,dashed,"v"]
\end{tikzcd}}

pour factoriser $u$ en
\[
u=vw
\]
où $w$ est une \uline{immersion fermée},
et $v$ est déduit de $\mathrm{pr}_1:I\times\mathord{\tikz[baseline=0pt,x=1em,y=1em]{\draw[line width=.045em,line join=round] (0,.02)--(.42,.68) (.17,.02)--(.52,.62)--(.83,.02)--cycle;}}^1$ par
\unsure{push-out|poush-out}. C'est une rétraction du
mapping-cone $C(u)$ sur l'ouvert $I'$, dont
le complémentaire fermé s'identifie à $I$ via $w$.
\scanpage{74}\mspage{61}
% archive sheet number: 73

Le foncteur de recollement entre $I',I$
\[
H:I'^{\unsure{o|\circ}}\times I\unsure{\text{---}|\to}\mathrm{Ens}
\]
est donné par
\[
\begin{aligned}
H(i',i)&=\mathrm{Hom}_{I'}(i',u(i))\\[-2pt]
\shortparallel\quad&\\[-2pt]
\inserted{\mathrm{Hom}_{C(u)}(i',i)}&
\end{aligned}
\tag*{$(*)$}
\]
La rétraction $v$ de $C(u)$ dans $I'$ applique
$i$ dans $u(i)$, et une flèche $i'\mathrel{\text{---}}i$
correspondant à $i'\longrightarrow u(i)$ dans $I'$ \inserted{via $(*)$}, est transformée
par $v$ en \unsure{cette} même flèche $i'\mathrel{\text{---}}u(i)$.

Calculons, pour $i'_0$ choisi, la catégorie
$\mathord{}_{i'_0}\backslash(C(u),v:C(u)\mathrel{\text{---}}I')$, des couples d'un
objet $\unsure{\zeta|z}$ de $C(u)$ et d'une flèche
$i'_0\longrightarrow v(\zeta)$ dans $I'$. On voit que $\mathord{}_{i'_0}\backslash C(u)$
admet l'ouvert $\mathord{}_{i'_0}\backslash I'$, avec
le complémentaire fermé $\mathord{}_{i'_0}\backslash I$,

% left margin sketch
\margin{\diagram{\begin{tikzcd}[ampersand replacement=\&,column sep=.1em,row sep=.5em,nodes={font=\tiny,inner sep=.1em}]
\underset{\mathrm{ouv.}}{I'} \&[-.6em] \subset \&[-.6em] C(u) \arrow[d] \&[-.6em] \unsure{\supset|\to} \&[-.6em] \underset{\substack{\text{fermé compl.}\\\text{de }I'}}{I}\\
\& \& I' \& \& 
\end{tikzcd}}}

ces deux catégories
étant recollées par le bifoncteur
\[
\begin{aligned}
(\mathord{}_{i'_0}\backslash I')^o\times \mathord{}_{i'_0}\backslash I&\longrightarrow\mathrm{Ens}\\
\big((\unsure{\alpha|\alpha_0}:\deleted{\unsure{\alpha}}i'_0\mathrel{\text{---}}i'),(\beta:i'_0\to\overset{\overset{i}{\vdots}}{u(i)})\big)
&\longmapsto\mathrm{Hom}_{\mathord{}_{i'_0}\backslash I'}(i',u(i))
\end{aligned}
\]
induit par $(*)$ ; i.e. les flèches de
$(\alpha:i'_0\mathrel{\text{---}}i')\unsure{\in}\mathrm{Ob}(\mathord{}_{i'_0}\backslash I')$ dans $(\beta:i'_0\unsure{\mapsto|\to|\text{---}}u(i))\in\mathord{}_{i'_0}\backslash I$
sont les flèches $i'\overset{\gamma}{\text{---}}u(i)$ telles que le
triangle

\diagram{\begin{tikzcd}[ampersand replacement=\&,row sep=small,column sep=small]
\& i'_0 \arrow[dl,"\alpha"'] \arrow[dr,"\beta"] \&\\
i' \arrow[rr,"\gamma"] \& \& u(i)
\end{tikzcd}}

commute.
\scanpage{75}\mspage{62}
% archive sheet number: 74

Si $\unsure{i'|a'|\alpha'}=i'_0$, $\alpha=\mathrm{id}_{i'_0}$, alors cette commutativité
revient à $\gamma=\beta$, \unsure{ça|ce} qui signifie qu'il
y a une flèche et une seule de
$\alpha$ dans $(i,\beta)$. Cela montre donc que
$(i'_0,\mathrm{id}_{i'_0})$ est un objet initial de $\mathord{}_{i'_0}\backslash C(u)$.
On a donc prouvé le lemme :

\uline{Lemme} 3.8. Soit \unsure{--} $u:I\unsure{\mapsto|\to} I'$ \deleted{\unsure{deux catégories}}\inserted{une flèche}
dans $\mathrm{Cat}$, factorisons-le en
\[
I\xrightarrow{w} C(u)\xrightarrow{v}I'\,,
\]
où $C(u)$ est le mapping-cone défini
p.ge 60, $w$ une immersion fermée
de complémentaire $I'$, $v$ la rétraction
canonique de $C(u)$ sur $I'$. Ceci
posé, pour \unsure{tt} $i'_0\in\mathrm{Ob}\, I'$, la catégorie
$\mathord{}_{i'_0}\backslash(C(\unsure{u|ui}),v)$ admet un objet initial ---
donc $v:C(u)\longrightarrow I'$ est (\unsure{tout ce|tant est} qu'il
y a de) \unsure{cofiltrante|cofiltrant}.

Nous pouvons maintenant prouver
3.7. Tout d'abord, les hypothèses
faites sur $I,I'$ se transportent à
$C(u)$ (hyp. artinienne, et \unsure{éventuelle|éventuelles}
hyp. \unsure{noethérienne|noethériennes} sur $\mathrm{Esp}(C(u))$. Donc
\scanpage{76}\mspage{63}
% archive sheet number: 75
% Paragraph continues from the preceding sheet.
on peut appliquer les résultats précédents,
notamment 3.5., aux flèches $v,w$. \deleted{\unsure{Le}}
\unsure{Lemme|Le lemme} précédent montre bien que $v$
satisfait aux conditions de 3.5, donc
le foncteur $v_!^W$ construit ds 3.3.1. est
adjoint à gauche de $v_W^*$. \inserted{Considérons} \deleted{\unsure{D}}d'autre part,
le foncteur $w_!^W$ construit \unsure{ds|dans} (3.3.1), je
dis que la $\unsure{=}$ conclusion s'applique.

Cela provient du fait que $w$ est une
immersion fermée. En effet, dans ce
cas il est immédiat que \deleted{$w_!:\unsure{u}$}
\[
w_!:M(I^{\unsure{o|\circ}})\unsure{\text{---}|\to} M(\deleted{\unsure{I}}J^o)\qquad(\text{où }J=C(u))
\]
est compatible avec les localiseurs,
puisque \unsure{plus|fibre} \unsure{à|a}
\[
\left\{\begin{aligned}
 w_!(X)(j)&\simeq X(j) &&\text{si }\deleted{\unsure{i}}j\in I\\
 w_!(X)(j)&\simeq\varnothing_M &&\text{si }j\in J\setminus I\,,
\end{aligned}\right.
\]
par le calcul des fibres.
\[
w_!(X)(j)\simeq\underrightarrow{\lim}_{i\in I/j}X(i)\,.
\]

Donc \unsure{le|la} \deleted{calcul}\inserted{définition} de $w_!^W$ \unsure{donné} dans
3.3.1. revient en fait à passer directement
aux cat. localisées pour $w_!$, sans avoir
à se restreindre à $M(I^o)_c$, et le
lemme des foncteurs adjoints nous dit que
\scanpage{77}\mspage{\unsure{64}}
% archive sheet number: 76
% Paragraph continues from the preceding sheet.
ce foncteur est adj. à g. de $w_W^*$ !

Comme on a évidemment
\[
u_!^W\simeq v_!^W w_!^W,\qquad u_W^*\simeq \deleted{\unsure{W}}w_W^*\circ v_W^{\unsure{'|*}}
\]
on en conclut que $u_!^W$ est bien adjoint
à gauche de $u_W^*$, \unsure{q ed.|q.e.d.}

\uline{Corollaire} 3.9. Soit $(M,W,C)$ une catégorie
à cofibrations. Alors le prédérivateur
$\mathbb{D}_W$, \deleted{\unsure{satisf}} restreint à la catégorie \deleted{\unsure{des}}
$\mathrm{Diag}_0$ formée des cat. \inserted{$I$} \uline{ordonnées} \deleted{\unsure{$I$}}
artiniennes et telles que $\mathrm{Esp}\,I$
soit noethérien, satisfait aux
axiomes Der 1 (\uline{décomposition}), Der 2
(\uline{\unsure{localisation|localisabilité}}), Der 3 (\uline{objet final}),
Der 4 bis (\uline{existence des $u_!$}), Der 5 bis
(mode de \uline{calcul des $u_!$}), Der 6
(\uline{exactitude des} carrés \inserted{$W$-}\uline{cocartésiens}).

Si $C$, $TC=C\cap W$ sont stables par composition
transfinie, on a le $\unsure{=}$ résultat
sur la catégorie $\mathrm{Diag}$ formée des
cat $I$ ordonnées artiniennes sans plus.
\scanpage{78}\mspage{65}
% archive sheet number: 77

3.10. Ainsi, pour avoir un dérivateur au plein
sens du terme, soit sur $\mathrm{Diag}_0$, soit sur
$\mathrm{Diag}$, il ne manque que Der 4 (existence
des $u_*$), Der 5 (calcul des $u_*$) et
Der 6 (exactitude des carrés $W$\deleted{\unsure{co}}-cartésiens).
On va \deleted{\unsure{l'}}\inserted{les} obtenir en principe à l'aide
d'une structure de \uline{catégorie à fibrations}
$(M,W,F)$ associée à $(M,W)$, \deleted{\unsure{ce}} i.e.
l'espèce de structure duale de celle
introduite \deleted{\unsure{dans}} p. 39. Ces \inserted{deux} structures
n'ont pas à être reliées entre elles d'une
façon \unsure{étroite|autre} que par le fait qu'elles
correspondent au $\unsure{=}$ $W$. Elles sont
considérablement plus faciles à \deleted{\unsure{dégag}}
construire que des structures de Quillen
$(W,C,F)$, où $C,F$ doivent être reliés
de façon \unsure{autrement|extrêmement} stricte, \deleted{\unsure{savoir}} \unsure{aux|par}
\deleted{\unsure{des prop}} les propriétés $C\longleftrightarrow\overset{F\cap W}{\overset{\shortparallel}{TF}}$, $\overset{C\cap W}{\overset{\shortparallel}{TC}}\unsure{\longleftarrow|\longleftrightarrow} F$,
(dont il n'est pas question dans l'axiomatique
de K.S. Brown), et \deleted{\unsure{de}}factorisation pour
$\unsure{\text{--}|(}C,TF)$ et $(TC,F)$.
\scanpage{79}\mspage{66}
% archive sheet number: 78

Mais \deleted{il reste} revenant à la construction
des $u_!^W$, il reste le pb de les
construire pour une flèche $u:I'\unsure{\text{---}|\to}I$
\uline{quelconque} dans $\mathrm{Cat}$.

À vrai dire, la construction du foncteur
\[
u_!^W:\mathrm{Ho}_W(I^{\unsure{o|\circ}})\mathrel{\text{---}}\mathrm{Ho}_W(I'^o)
\tag*{(3.3.1)}
\]
est valable sans supposer que $I'$ soit une
catégorie ordonnée --- \unsure{ça} peut être une
catégorie quelconque. Cela provient du
fait que le foncteur induit par $u_!$
\[
u_{!\,\mathrm{cof}}:M(I^o)_c\longrightarrow M(I'^o)
\]
\unsure{soit|doit} encore compatible avec les localiseurs,
comme on voit par le
calcul des fibres
\[
u_!(\unsure{x|X})(i')\simeq \underrightarrow{\lim}_{i\in I/\unsure{i'|u'}}X(i)\,.
\]

Ainsi, on est ramené au cas où $I'=e$,
quitte à remplacer $I$ par $I/i'$, et en
notant que ($I/\unsure{j'|i'}\mathrel{\text{---}}I$ étant étale) l'image
inverse de $X$ sur $I/i'$ est \unsure{cofibrant|cofibrante}
(cf. 3.2. a).
\scanpage{80}\mspage{67}
% archive sheet number: 79

Mais il \unsure{se|pr} présente à nouveau la perplexité
si le foncteur (3.3.1) ainsi construit
est bien un adjoint à gauche de $u_W^*$.
On peut songer à utiliser encore la
$\unsure{=}$ factorisation en $u=vw$, en utilisant
le mapping-cône. La difficulté, cependant,
c'est qu'à présent $J=C(u)$ n'est plus une
cat ordonnée, donc on n'a pas une
définition toute prête de \deleted{\illegible{1}} $C_{J^{\unsure{o|\circ}}}$
$\subset\mathrm{Fl}\, M(J^o)$, donnant lieu à
une sous-catégorie $M(\deleted{\unsure{\in}}J^o)_c$ etc. En
fait, il semblerait qu'on soit pratiquement
paralysé, \unsure{tant} qu'on n'arrive à construire
sur $M(J^o)$ également une structure
de catégorie à \unsure{cofibration|cofibrations}, se comportant
bien par rapport au foncteur $\unsure{v|w}_!$.
La ``démonstration'' d'Anderson de l'existence
de $u_!$ pour \uline{\unsure{tt}} flèches dans $\mathrm{Cat}$,
foire totalement, \deleted{\unsure{car}} ($=$ pour $I\unsure{\text{---}|\to}e$,
$I$ un ordonné fini --- mais là j'ai
pu rétablir la situation), car il se ramène
au cas où $I$ est ordonné loc. fini, et il
\scanpage{81}\mspage{68}
% archive sheet number: 80
% Paragraph continues from the preceding sheet.
admet sans plus que \uline{dans ce cas} $u_!^W$ est
adjoint à gauche de $u_W^*$ ; et \uline{de plus},
que si $u$ est $W_0$-\deleted{\unsure{homotopique}}\inserted{\unsure{asphérique}} donc
$u^*:M(I'^{\unsure{o|\circ}})\unsure{\text{---}|\to}M(I^o)$ pl. fidèle, qu'il en
soit de même de $u_W^{\unsure{'|1|*}}$ déduit par passage
aux localisés. Or même si on admet
que $u_!^W$ est adj. à gauche de $u_{\unsure{*|\times}}^W$,
faute de disposer d'un morphisme
d'adjonction tout prêt ``en bas'' (dans
$\mathrm{Ho}_W(I'^o)$), déduit du morphisme
d'adjonction dans $M(I'^o)$, il n'est
absolument pas clair que \deleted{\unsure{$u_!^W$}}\unsure{ce}
$u_W^*$ soit pl. fid. i.e. $u_!^W$ un foncteur
de localisation. Le contre-exemple 3.6.
montre bien que, $\unsure{=}$ si on a adjonction
(i.e. si $e$ est un objet initial de $C=C'$],
le foncteur $u^*=\mathrm{id}_C$ a beau être pl.
fid (et $\mathord{=}$ une équivalence de catégories),
le foncteur localisé n'a aucune raison
de l'être : cela prouverait que \unsure{tt} objet
\deleted{catégorie}\inserted{de $\mathrm{Cat}$} ayant un objet initial est
\deleted{un groupoïde} équivalent à $e$ !
\scanpage{82}\mspage{69}
% archive sheet number: 81

3.11. \unsure{Disons|Dirons} qu'on a une \uline{catégorie à \unsure{cofibration|cofibrations}
stricte} si dans l'axiome $\mathrm{Cof}\,3$ de factorisation
(cf p. 39), la factorisation peut être prise
\deleted{\uline{\unsure{raisonnablement}}}\inserted{fonctoriellement}. Il revient \inserted{$\pm$} au $\unsure{=}$ de dire que
l'axiome de factorisation reste valable,
quand on passe à une catégorie $M(I)$
pour $I$ dans $\mathrm{Cat}$, en y définissant $W(I)$ et
$C(I)$ (à ne pas confondre avec $C_I$) fibre par
fibre. Cela revient au $\mathord{=}$ \deleted{\unsure{tout}} au moins si
$M$ est équiv. à une \unsure{petite} catégorie. Sinon,
il faut supposer fonctorialité pour toute
\uline{petite} sous-catégorie pleine de $\uline{\mathrm{Fl}}\,M$.
Il vaut mieux définir la catégorie à cofibrations
stricte par cette propriété de stabilité.

On voit alors que le prédérivateur
$\mathbb{D}=\mathbb{D}_{\unsure{M|W}}$ \deleted{\unsure{à la}} sur $\mathrm{Cat}$ est \deleted{\illegible{1}} un dérivateur
\uline{relativement à $D=\mathrm{Diag}_0$} défini plus haut
(cat. ordonnées \inserted{$I$} artiniennes, telles que $\mathrm{Esp}(I)$
soit noethérien), ou $\mathord{=}$ rel. à \inserted{$D=$}$\mathrm{Diag}$ (cat
ordonnées $I$ artiniennes sans plus)
lorsque $C$ et $TC$ sont stables par composition
transfinie. J'entends par là que
$u_!$ existe quand $u$ est une flèche de $D$,
\scanpage{83}\mspage{70}
% archive sheet number: 82
% Paragraph continues from the preceding sheet.
et que ça reste valable quand on passe
à un \unsure{dérivateur} induit $\mathbb{D}_Z$ ($Z$ une \unsure{cat|Cat}
quelc.) ; en d'autres termes, que $\unsure{v|u}_!$ existe
quand $v$ est de la forme $\mathrm{id}_Z\times u$, où $u:I\to I'$
est une flèche de $D$, et que le
calcul se fait \uline{au dessus de $I'$} selon
la règle Dér 5 bis. Cela signifie donc que
la flèche de changement de base pour les
$f_!$ est un iso dans les cas suivants :

\diagram{\begin{tikzcd}[ampersand replacement=\&,column sep=huge,row sep=large]
I\times Z \arrow[d,"u\times\mathrm{id}_Z"] \& \mathord{}_{i'}\backslash I\times Z \arrow[l] \arrow[d,"\mathrm{pr}_2"]\\
I'\times Z \& \deleted{\unsure{\mathord{}_{i'}\backslash I}}\{i'\}\times Z\,, \arrow[l]
\end{tikzcd}}

où $i'$ est un objet quelc. de $I'$.

L'hypothèse \uline{stricte} pour une catégorie à
cofibrations me \unsure{parait|paraît} assez \unsure{anodine} --- dans
la plupart des cas \unsure{à|:} ma connaissance, elle
est satisfaite. Mais \deleted{\unsure{de ma chose}} $\unsure{=}$
moyennant cette hypothèse, je ne vois
aucun moyen de prouver l'existence
des $u_!$ pour \uline{\unsure{tt}} flèches de $\mathrm{Cat}$, et c'est
une situation bien sûr très gênante.
On voudrait \unsure{absolument|isolément|seulement} dégager des
conditions qui assurent la validité de Der 4 bis
et de Der 5 bis (existence des $u_!$ pour
\scanpage{84}\mspage{71}
% archive sheet number: 83
% Paragraph continues from the preceding sheet.
\uline{toutes} flèches de $\mathrm{Cat}$, et mode de calcul des
fibres \ldots).

4. \uline{Cat. à cofibrations spéciales}

4.1. Je suppose à présent que \textcircled{a} $M$ est stable par
$\underrightarrow{\lim}$ filtrantes, et \textcircled{b} que ses objets sont
accessibles, de sorte qu'on est sous les
conditions d'application du th. de factorisation
(\unsure{X} \S\deleted{\unsure{\S}}1).
% slanted
\margin{\textcircled{c} donnés \inserted{$W$,} $C\subset\mathrm{Fl}\,M$ satisfaisant les conditions Cof 1, Cof 2 de \unsure{p.} 39. Pour Cof 3, cf plus bas.}
\deleted{\unsure{Pos}} On suppose \textcircled{d} $C,TC$\inserted{$=C\cap W$} \uline{stables \unsure{par|pour}
composition transfinie} (c'est le cas si la
structure provient d'une structure de
Quillen close $(W,C,F)$).
\deleted{On suppose de
plus que \textcircled{d}
\[
TF\overset{\mathrm{déf}}{=}C_*\subset W\tag*{(4.1)}
\]
donc $TF\subset W_{\mathrm{univ}}$ (puisqu'il est stable par
chgt. de base). \deleted{\unsure{ce qui est|de surcroît} \unsure{stable|stables} \unsure{on suppose} \unsure{que} \illegible{1}}}
\deleted{\unsure{si on désire, donc}} Posons
\[
\deleted{\unsure{=}}\quad TF\overset{\mathrm{déf}}{=}C_*\subset\mathrm{Fl}\,M\,,
\tag*{(4.1)}
\]
il résulte du th. de factorisation que \deleted{\unsure{toute}}\unsure{tte}
flèche $f$ de $M$ se factorise en
\[
f=pi,\qquad p\in TF,\ i\in C
\tag*{(4.1)}
\]
\uline{pourvu} qu'on fasse l'hypothèse \textcircled{\unsure{e|d}}
\[
\begin{aligned}
&\exists\, C_0\subset C,\quad \text{\uline{petite} partie de }C\text{, telle}\\
&\text{que}\quad C\subset\widetilde{C_0}\quad \text{i.e.}\quad TF=(C_0)_*\,.
\end{aligned}
\tag*{(4.3)}
\]
Sous ces conditions, si on suppose \textcircled{\unsure{f|g}}
\[
TF\subset W\,,
\tag*{(4.4.)}
\]
